% from ega4-2-fr.tex
\oldpage[IV-2]{101}

On a $\dim(Z_i)\leq\dim(Y)$ pour tout $i$; d'autre part, si $\zeta_i$ est le
point générique de $Z_i$, on a $\dim(g_i^{-1}(\zeta_i))\leq n$ pour tout $i$
par hypothèse; on voit ainsi que pour prouver (5.6.7.1), il suffit de
démontrer que $\dim(X_i)\leq\dim(Z_i)+n$ pour tout $i$, ce qui résulte de
(5.6.6.1).

\begin{corollary}[5.6.8]
\label{IV.2.s5.6.8-fr}
Soient $Y$ un préschéma localement noethérien, $X$ un préschéma
irréductible, $f:X\to Y$ un morphisme localement de type fini, $n$ un entier
$\geq0$. On suppose que $Y$ est universellement caténaire, et que pour tout
$y\in Y$, on ait $\dim(f^{-1}(y))\geq n$ (resp.
$\dim(f^{-1}(y))=n$). Alors on a
\begin{equation}
\dim(X)\geq\dim(Y)+n
\tag{5.6.8.1}\label{IV.2.s5.6.8.1-fr}
\end{equation}
(resp.
\begin{equation}
\dim(X)=\dim(Y)+n).
\tag{5.6.8.2}\label{IV.2.s5.6.8.2-fr}
\end{equation}
\end{corollary}

Comme $n\geq0$, l'hypothèse entraîne que $f$ est surjectif, donc $Y$
irréductible; en outre, si $\eta$ est le point générique de $Y$, on a
$\dim(f^{-1}(\eta))\geq n$ (resp. $\dim(f^{-1}(\eta))=n$). La conclusion
résulte donc de (5.6.6).

\begin{remarks}[5.6.9]
\label{IV.2.s5.6.9-fr}
\begin{enumerate}
\item[(i)] Même si $Y$ est régulier, $X$ irréductible, $f:X\to Y$ dominant
et de type fini, les deux membres de (5.6.6.1) ne sont pas nécessairement
égaux, comme le montre l'exemple où $Y=\operatorname{Spec}(A)$ où $A$ est un
anneau de valuation discrète, $X=\operatorname{Spec}(K)$ où $K$ est le corps
des fractions de $A$, $f$ étant le morphisme canonique.
\item[(ii)] L'exemple (5.4.3, (i)) montre que dans (5.6.6) et (5.6.8), on ne
peut supprimer l'hypothèse que $X$ est irréductible, les autres hypothèses
étant vérifiées. Nous verrons toutefois (10.6.1) que moyennant des hypothèses
supplémentaires sur $Y$, vérifiées par exemple lorsque $Y$ est un préschéma
localement de type fini sur un corps, ou sur un anneau de Dedekind ayant une
infinité d'idéaux premiers, de tels phénomènes ne peuvent se présenter.
\end{enumerate}
\end{remarks}

\begin{proposition}[5.6.10]
\label{IV.2.s5.6.10-fr}
Soient $A$ un anneau local noethérien intègre universellement caténaire, $B$
un anneau intègre contenant $A$ et qui est une $A$-algèbre finie. Alors, pour
tout idéal maximal $\mathfrak n$ de $B$, on a
$\dim(B_{\mathfrak n})=\dim(A)$.
\end{proposition}

En effet, on a $\operatorname{deg.tr.}_A B=0$ et le corps résiduel $k'$ de
$B_{\mathfrak n}$ est une extension algébrique du corps des fractions de
$A$. La conclusion résulte de la formule (5.6.1.1), tout idéal maximal de
$B$ étant au-dessus de celui de $A$.

\begin{example}[5.6.11]
\label{IV.2.s5.6.11-fr}
Soit $A$ un anneau local noethérien intègre et intégralement clos; alors, on
a vu (0, 16.1.6) que la conclusion de (5.6.10) est valable pour toute
$A$-algèbre finie intègre $B$ contenant $A$. Par contre, nous allons
construire un exemple d'anneau local noethérien intègre caténaire $A$ pour
lequel la conclusion de (5.6.10) sera en défaut. Nous utiliserons la
construction de (5.2.5, (i)). Soient $k_0$ un corps, $k$ une extension
transcendante pure $k_0(S_i)_{i\in\mathbb N}$ de degré de transcendance
infini, $V$ l'anneau de valuation discrète $k[S]_{(S)}$, anneau local de
l'anneau de polynômes $k[S]$ en l'idéal premier $(S)$; enfin soit $E$
l'anneau de polynômes $V[T]$; on a vu que dans $E$ l'idéal maximal
$\mathfrak m=(S)+(T)$ est de hauteur 2 et l'idéal maximal
$\mathfrak m'=(ST-1)$ de hauteur 1;

